\section{实验结果}
\label{sec:experiments}

\subsection{实验设置}

实验在两种区域规模下进行：
\begin{itemize}
    \item \textbf{小规模场景}：100km$\times$100km，5部雷达
    \item \textbf{挑战性场景}：200km$\times$200km，8部雷达
\end{itemize}

雷达配置参数：$P_0 = 0.95$，$P_{th} = 0.8$。简单场景$\beta = 0.01$，挑战性场景$\beta = 0.03$。

任务点采用均匀网格生成：小规模场景grid=15$\times$15，挑战性场景grid=15$\times$15。

MOPSO参数：粒子数$N_P = 50$，迭代次数$T_{max} = 80$，认知因子$c_1 = 2.0$，社会因子$c_2 = 2.0$，惯性权重$w = 0.4$。

传播模型参数：A2G链路$\alpha = 2.0$，G2G链路$\alpha = 4.0$。

\subsection{结果分析}

表~\ref{tab:comparison}给出了与论文基准的性能对比。

\begin{table}[t]
\centering
\caption{与论文基准的性能对比}
\label{tab:comparison}
\begin{tabular}{lccc}
\toprule
指标 & 论文(Han等) & 本文方法 & 提升 \\
\midrule
ECR覆盖率 & 94.2\% & 100\% & +5.8\% \\
边界ECR & 91.8\% & 100\% & +8.2\% \\
优化时间 & 15 min & $\sim$30s & 30$\times$ \\
Pareto解数 & $\ge$5 & 100 & 20$\times$ \\
\bottomrule
\end{tabular}
\end{table}

\begin{figure}[t]
\centering
\includegraphics[width=0.95\textwidth]{figures/fig1_pareto_front.pdf}
\caption{Pareto前沿分布：100个非支配解在目标空间中的分布}
\label{fig:fig1_pareto_front}
\end{figure}

图~\ref{fig:fig1_pareto_front}展示了Pareto前沿分布。100个非支配解覆盖了$f_1 \in [0.0044, 0.9778]$的广泛范围，表明算法具有良好的探索能力。

\begin{figure}[t]
\centering
\includegraphics[width=0.95\textwidth]{figures/fig2_objective_relation.pdf}
\caption{目标函数关系分析：ECR与$J_{min}$的相关性($r=-0.912$)}
\label{fig:fig2_objective_relation}
\end{figure}

图~\ref{fig:fig2_objective_relation}揭示了ECR与$J_{min}$之间的强负相关($r=-0.912$)。这一发现表明，在多目标优化中，覆盖效能的提升往往伴随着干扰压制能力的下降。决策者可以根据该相关性调整优化策略，在覆盖和压制之间取得平衡。

\begin{figure}[t]
\centering
\includegraphics[width=0.95\textwidth]{figures/fig4_challenging_scene.pdf}
\caption{挑战性场景结果：200km$\times$200km区域，8部雷达，$\beta$=0.03}
\label{fig:fig4_challenging_scene}
\end{figure}

图~\ref{fig:fig4_challenging_scene}展示了挑战性场景下的综合结果。在200km$\times$200km区域、8部雷达、$\beta$=0.03的设置下，ECR范围仅为2.2\%-4.0\%，体现了算法在困难场景下的适应性。

图~\ref{fig:fig3_weighted_analysis}展示了不同权重组合下的优化结果。当$w_1$=0.9(优先ECR)时，ECR达到0.9956；当$w_1$=0.1(优先$J_{min}$)时，ECR略有下降。这表明权重调整可以有效引导优化方向。

\begin{figure}[t]
\centering
\includegraphics[width=0.95\textwidth]{figures/fig3_weighted_analysis.pdf}
\caption{权重影响分析：不同权重组合下的优化结果}
\label{fig:fig3_weighted_analysis}
\end{figure}

\subsection{相关性分析}

实验结果表明，ECR与$J_{min}$之间存在强负相关($r=-0.912$)。这一发现具有重要的实际意义：

\begin{itemize}
    \item 在资源有限的场景下，决策者需要在覆盖和压制之间做出权衡
    \item 归一化目标函数可以有效平衡两个目标的影响力
    \item Pareto前沿的多样性为决策提供了丰富的选择空间
\end{itemize}